\section{Related work}

\paragraph{NLA and mechanistic interpretability.}
A long line of interpretability work reads structure out of model internals, from
linear probes that decode features from intermediate
representations~\cite{alain2017probing} to sparse-dictionary methods that recover
interpretable directions in superposition~\cite{bricken2023monosemanticity,
templeton2024scaling}. Natural Language Autoencoders~\cite{nla2026} push this
further by training a decoder that verbalizes an activation directly into natural
language, so that a model's internal state can be \emph{read} as text. Our use of
the NLA idea is narrow and instrumental: we treat the verbalizer's text, scored by
a classifier, as the \emph{definition} of a concept, which lets a downstream
component decide ``does this internal state read as test-awareness?'' without our
having to hand-pick an activation direction. We stress that we do not have a trained
NLA locally; our verbalizer is a retrieval-plus-projection stand-in
(Section~\ref{sec:method}), and every readout it emits is tagged as such.

\paragraph{Certified robustness and convex relaxations.}
Empirical adversarial robustness can always be undercut by a stronger attack, which
motivated \emph{certified} defenses that produce a one-sided guarantee. Raghunathan,
Steinhardt and Liang gave a semidefinite and a dual/linear relaxation that upper-bound
the worst-case margin of a one-hidden-layer ReLU network over an $L_\infty$
ball, and --- the part we lean on most --- made that bound differentiable so it can
be optimized as a regularizer~\cite{raghunathan2018certified,raghunathan2018sdp}.
Parallel and follow-on work includes the convex outer adversarial
polytope~\cite{wong2018provable}, CROWN's general-activation backward linear
relaxation~\cite{zhang2018crown}, interval bound propagation~\cite{gowal2018ibp},
and their combination CROWN-IBP~\cite{zhang2020crownibp}, later packaged in the
\texttt{auto\_LiRPA} library~\cite{xu2020autolirpa}. We use this machinery in an
unusual direction. Standard certified defense asks ``can an adversary flip the
predicted label inside a ball?''; we ask ``can \emph{any} input inside a ball reach
a designated internal region $A^\star$?'', which is the same worst-case-output
computation read as a reachability question. Because $A^\star$ is the decision
region of a one-hidden-layer ReLU network by construction, the 2018 certificate
applies essentially verbatim.

\paragraph{Discrete inversion.}
On the constructive side, finding a real prompt that lands in $A^\star$ is a
discrete search over tokens. We use HotFlip's gradient linearization of the
one-hot input~\cite{ebrahimi2018hotflip} inside a greedy coordinate-gradient (GCG)
loop~\cite{zou2023gcg}, scoring candidate swaps against the true model and surrogate.
This is the same family of methods used to attack aligned LLMs, here aimed at
crossing a concept threshold rather than eliciting a harmful string.

\paragraph{Evaluation-awareness.}
The phenomenon we target is now documented in frontier-model reporting. For Claude
Sonnet~4.5, evaluation-awareness was verbalized in roughly $16\%$ of code-safety
test transcripts and roughly $26\%$ of SWE-bench-style transcripts, versus under
$1\%$ of real user traffic, and a political-sycophancy test was explicitly flagged
in about $13\%$ of transcripts~\cite{anthropic2025sonnet45}. These numbers are our
motivation, not our data: our experiments run on GPT-2 with a synthetic
eval-vs-deploy task, and the realistic rates above describe the gap we ultimately
want a certified method to speak to.

\paragraph{What a certificate does \emph{not} answer.}
A reachability certificate is about \emph{inference-time} inputs: which prompts can
or cannot drive the model into $A^\star$. It says nothing about which
\emph{training} examples instilled the concept in the first place. That second
question --- attributing a behavior back to data --- is the province of influence
functions and related attribution methods~\cite{koh2017influence,grosse2023influence},
and is out of scope here. We make this distinction explicit because an early framing
of this project conflated the two, and the correction matters for how the result is
read (Section~\ref{sec:limitations}).
